\documentclass[a4paper,11pt,reqno]{amsart}

\usepackage[left=1in,right=1in,top=1in,bottom=1in]{geometry}
\usepackage{enumitem}
\usepackage{amssymb,amsmath,latexsym,amsfonts,amsbsy,amsthm,mathtools,color}
\usepackage{float}
\usepackage[colorlinks=true,citecolor=blue,filecolor=blue,linkcolor=blue,urlcolor=blue]{hyperref}

% Load subject-specific packages only when the manuscript uses them.

% Canonical personal macro library, adapted from the bundled arXiv v2
% template example. These definitions are protected by deterministic checks.
\newcommand{\al}{\alpha}
\newcommand{\vt}{\vartheta}
\newcommand{\lda}{\lambda}
\newcommand{\om}{\Omega}
\newcommand{\pa}{\partial}
\newcommand{\va}{\varepsilon}
\newcommand{\ud}{\mathrm{d}}
\newcommand{\be}{\begin{equation}}
\newcommand{\ee}{\end{equation}}
\newcommand{\w}{\omega}
\newcommand{\Lda}{\Lambda}

\newcommand{\A}{\mathbf{A}}
\newcommand{\cA}{\mathcal{A}}
\newcommand{\bA}{\mathbb{A}}
\newcommand{\B}{\mathbf{B}}
\newcommand{\bb}{\mathbf{b}}
\newcommand{\cB}{\mathcal{B}}
\newcommand{\CC}{\mathbf{C}}
\newcommand{\cC}{\mathcal{C}}
\newcommand{\C}{\mathbb{C}}
\newcommand{\bd}{\mathbf{d}}
\newcommand{\bD}{\mathbb{D}}
\newcommand{\rD}{\mathrm{D}}
\newcommand{\cD}{\mathcal{D}}
\newcommand{\E}{\mathbb{E}}
\newcommand{\cE}{\mathcal{E}}
\newcommand{\e}{\mathbf{e}}
\newcommand{\re}{\mathrm{e}}
\newcommand{\fF}{\mathbf{F}}
\newcommand{\cF}{\mathcal{F}}
\newcommand{\bG}{\mathbb{G}}
\newcommand{\fG}{\mathbf{G}}
\newcommand{\cG}{\mathcal{G}}
\newcommand{\rh}{\mathrm{h}}
\newcommand{\rH}{\mathrm{H}}
\newcommand{\bH}{\mathbb{H}}
\newcommand{\I}{\mathbf{I}}
\newcommand{\ii}{\mathrm{i}}
\newcommand{\jj}{\mathrm{j}}
\newcommand{\kk}{\mathrm{k}}
\newcommand{\LL}{\mathbf{L}}
\newcommand{\cL}{\mathcal{L}}
\newcommand{\Z}{\mathbb{Z}}
\newcommand{\M}{\mathcal{M}}
\newcommand{\MM}{\mathbb{M}}
\newcommand{\bM}{\mathbf{M}}
\newcommand{\m}{\mathbf{m}}
\newcommand{\cN}{\mathcal{N}}
\newcommand{\N}{\mathbb{N}}
\newcommand{\rN}{\mathrm{N}}
\newcommand{\n}{\mathbf{n}}
\newcommand{\OO}{\mathbb{O}}
\newcommand{\cP}{\mathcal{P}}
\newcommand{\PP}{\mathbf{P}}
\newcommand{\p}{\mathbf{p}}
\newcommand{\Q}{\mathbf{Q}}
\newcommand{\R}{\mathbb{R}}
\newcommand{\RR}{\mathbf{R}}
\newcommand{\cR}{\mathcal{R}}
\newcommand{\cS}{\mathcal{S}}
\newcommand{\Ss}{\mathbb{S}}
\newcommand{\fS}{\mathbf{S}}
\newcommand{\cT}{\mathcal{T}}
\newcommand{\uu}{\mathbf{u}}
\newcommand{\U}{\mathbf{U}}
\newcommand{\vv}{\mathbf{v}}
\newcommand{\V}{\mathbf{V}}
\newcommand{\W}{\mathbf{W}}
\newcommand{\cW}{\mathcal{W}}
\newcommand{\X}{\mathbf{X}}
\newcommand{\cX}{\mathcal{X}}
\newcommand{\Y}{\mathbf{Y}}
\newcommand{\cY}{\mathcal{Y}}

\newcommand{\wc}{\rightharpoonup}
\newcommand{\Sn}{\Ss^n}
\newcommand{\HH}{\mathcal{H}}
\newcommand{\FF}{\mathcal{F}}
\newcommand{\vp}{\varphi}
\newcommand{\da}{\downarrow}
\newcommand{\T}{\mathrm{T}}
\newcommand{\ga}{\gamma}
\newcommand{\Ga}{\Gamma}
\newcommand{\ep}{\epsilon}
\newcommand{\sg}{\sigma}
\newcommand{\ift}{\infty}
\newcommand{\wt}{\widetilde}
\newcommand{\wh}{\widehat}
\newcommand{\f}{\frac}
\newcommand{\D}{\mathrm{D}}
\newcommand{\ol}{\overline}
\newcommand{\Ra}{\Rightarrow}
\newcommand{\op}{\operatorname}
\newcommand{\Sg}{\Sigma}
\newcommand{\nn}{\nonumber}
\newcommand{\na}{\nabla}
\newcommand{\dint}{\displaystyle\int}
\newcommand{\dlim}{\displaystyle\lim}
\newcommand{\dsup}{\displaystyle\sup}
\newcommand{\dmin}{\displaystyle\min}
\newcommand{\dmax}{\displaystyle\max}
\newcommand{\dinf}{\displaystyle\inf}
\newcommand{\dsum}{\displaystyle\sum}

% Delimiter enlargement uses \left and \right only.
\newcommand{\abs}[1]{\left\lvert#1\right\rvert}
\newcommand{\norm}[1]{\left\lVert#1\right\rVert}
\newcommand{\inner}[2]{\left\langle#1,#2\right\rangle}
\newcommand{\paren}[1]{\left(#1\right)}
\newcommand{\dd}{\,\mathrm{d}}

\DeclareMathOperator{\dist}{dist}
\DeclareMathOperator{\avg}{avg}
\DeclareMathOperator{\Bad}{Bad}
\DeclareMathOperator{\diam}{diam}
\DeclareMathOperator{\BMO}{BMO}
\DeclareMathOperator{\VMO}{VMO}
\DeclareMathOperator{\sgn}{sgn}
\DeclareMathOperator{\supp}{supp}
\DeclareMathOperator{\im}{im}
\DeclareMathOperator{\Meas}{Meas}
\DeclareMathOperator{\pts}{pts}
\DeclareMathOperator{\Lip}{Lip}
\DeclareMathOperator{\Vol}{Vol}
\DeclareMathOperator{\tr}{tr}
\DeclareMathOperator*{\osc}{osc}
\DeclareMathOperator{\Length}{Length}
\DeclareMathOperator{\diag}{diag}
\DeclareMathOperator{\graph}{graph}
\DeclareMathOperator{\even}{even}
\DeclareMathOperator{\odd}{odd}
\DeclareMathOperator{\sing}{sing}
\DeclareMathOperator{\reg}{reg}
\DeclareMathOperator{\good}{good}
\DeclareMathOperator{\bad}{bad}
\DeclareMathOperator{\rank}{rank}
\DeclareMathOperator{\const}{const}
\DeclareMathOperator{\End}{End}
\DeclareMathOperator{\loc}{loc}
\DeclareMathOperator{\argmin}{argmin}
\DeclareMathOperator{\Id}{Id}
\DeclareMathOperator{\Rei}{Rei}
\DeclareMathOperator{\ad}{ad}
\DeclareMathOperator{\id}{id}
\DeclareMathOperator{\Int}{Int}
\DeclareMathOperator{\fun}{fun}
\DeclareMathOperator{\curl}{curl}
\DeclareMathOperator{\CN}{CN}
\DeclareMathOperator{\HSV}{HSV}
\DeclareMathOperator{\Center}{Center}
\DeclareMathOperator{\vol}{vol}
\DeclareMathOperator*{\esssup}{ess\,sup}
\DeclareMathOperator*{\essinf}{ess\,inf}

\setlist[enumerate,1]{label=$(\theenumi)$,ref=\theenumi,leftmargin=*}
\numberwithin{equation}{section}

\theoremstyle{plain}
\newtheorem{thm}{Theorem}[section]
\newtheorem{cor}[thm]{Corollary}
\newtheorem{con}[thm]{Conjecture}
\newtheorem{lem}[thm]{Lemma}
\newtheorem{prop}[thm]{Proposition}
\newtheorem{assum}[thm]{Assumption}

\theoremstyle{definition}
\newtheorem{defn}[thm]{Definition}
\newtheorem{example}[thm]{Example}
\newtheorem{q}[thm]{Question}

\theoremstyle{remark}
\newtheorem{rem}[thm]{Remark}

\title{A sharp inequality for two real numbers}
\author{Example Author}
\date{}
\begin{document}
\begin{abstract}
For real numbers, the square of their sum is bounded by twice the sum of their squares, with equality precisely when the two numbers coincide. The coefficient two is the smallest constant for which this estimate holds for every pair of real numbers. This elementary result requires no sign restriction on either variable and includes zero and opposite-sign inputs.
\end{abstract}
\subjclass[2020]{26D15}
\keywords{Elementary inequalities, sharp constants, equality cases}
\maketitle

\section{Introduction}

\subsection{Main results}
Let $a$ and $b$ be real numbers. An optimal coefficient is a constant valid for every pair of inputs that cannot be replaced by a smaller constant.
\begin{thm}\label{thm:sum}
For every pair $a,b\in\R$,
\[
 (a+b)^2\leq 2(a^2+b^2).
\]
Equality holds if and only if $a=b$, and the coefficient $2$ cannot be decreased.
\end{thm}

\subsection{Related results}
The classical Cauchy--Schwarz inequality for real Euclidean vectors has the form
\[
 \abs{\inner{v}{w}}^2\leq\norm{v}^2\norm{w}^2.
\]
Taking $v=(a,b)$ and $w=(1,1)$ recovers the present scalar inequality. This note studies that special case and makes no novelty claim.

\subsection{Difficulties and strategies}
The mixed term $2ab$ cannot be discarded merely because the square of the difference is non-negative. The useful comparison is between that mixed term and the sum of the two individual squares.

Subtracting the square of the sum from twice the sum of the squares isolates one non-negative square. Its vanishing identifies the equality case, while any non-zero equal pair tests the necessity of the coefficient.

\subsection{Organization of this paper}
Section~\ref{sec:preliminaries} supplies the elementary sign fact needed for the estimate. Section~\ref{sec:proof} establishes the inequality, equality condition, and optimal coefficient.

\section{Preliminaries}\label{sec:preliminaries}
The order on the real numbers controls the sign of a square.
\begin{lem}\label{lem:square}
For every $x\in\R$, one has $x^2\geq0$, with equality if and only if $x=0$.
\end{lem}
\begin{proof}
If $x>0$, multiplication by $x$ yields $x^2>0$.
\end{proof}

\section{The sharp estimate}\label{sec:proof}
The square of the difference controls the mixed term in the desired estimate.
\begin{proof}[Proof of Theorem~\ref{thm:sum}]
Lemma~\ref{lem:square} gives $(a-b)^2\geq0$. Consequently, $2ab\leq0$ and
\[
 (a+b)^2=a^2+b^2+2ab\leq a^2+b^2\leq 2(a^2+b^2).
\]
\end{proof}
\end{document}
